\chapter{Project management}
\label{chap:project-management}
Throughout the project, our team adopted a hybrid approach to project
management, combining traditional planning with Agile-inspired iterative
development. This strategy enabled us to maintain organization while
remaining responsive to evolving requirements and challenges.

\subsubsection{Planning and Tooling}\label{Planning and Tooling}

Initially, we used \textbf{Trello} to outline key deliverables and
divide the project into clear, actionable tasks. Each task was clearly
assigned and categorized into columns such as "To Do," "In Progress,"
and "Completed," enhancing transparency and accountability among team
members.

Version control and collaboration were handled through \textbf{GitHub},
where each team member created a new branch every week. This approach
allowed us to clearly manage version control by individual contribution
rather than by feature, making it easier to pinpoint issues and revert
changes if necessary. Pull requests were regularly reviewed and approved
by teammates before merging into our main development branch, ensuring
code quality and consistency.

\subsubsection{Iterative Progress and Agile Practices}\label{Iterative Progress and Agile Practices}

While we initially created a comprehensive roadmap, we consistently
applied Agile practices, holding weekly meetings to review progress and
adjust our priorities accordingly. For example, three weeks before the
final deadline, we discovered our goal of creating a wide variety of
customizable components in the Club Hub was overly ambitious given our
timeline. Consequently, we decided to simplify this feature,
implementing only essential components such as images, clocks, and text.
However, we strategically designed our codebase to facilitate future
expansions.

Another example occurred during the middle of our development when we
identified that implementing a user deletion feature was highly
problematic due to extensive foreign key dependencies in the database.
Recognizing the risk to data integrity, we proactively decided to
abandon this feature. Such flexibility significantly reduced potential
issues and maintained project stability.

\subsubsection{Communication and Collaboration}\label{Communication and Collaboration}

Our communication extended beyond scheduled meetings, as team members
frequently interacted privately through messaging platforms to address
tightly coupled features. For instance, frontend and backend developers
regularly coordinated to ensure seamless integration of the event
management functionalities, which required continuous adjustments to the
Django views and frontend templates. This constant dialogue was crucial
in resolving unexpected issues promptly and maintaining steady progress.

\subsubsection{Evaluation and Lessons Learned}\label{Evaluation and Lessons Learned}

In this project, we effectively utilized version control tools to ensure
that each team member had continuous access to the latest code
developments. We used Trello to track the tasks and progress of each
member effectively. However, during the initial phases, we inaccurately
assessed our project architecture, misunderstanding
Django\textquotesingle s framework features, and incorrectly allocated
tasks based on frontend and backend roles. Although we identified this
issue by the third week of the project and switched our task allocation
strategy to assigning tasks based on specific features, this initial
mistake inevitably caused tight coupling in our code. As a result,
despite consistent effort and contributions from all team members, we
experienced difficulties distributing tasks evenly in the
project\textquotesingle s later stages.
